\section{Aims and objectives \label{sec:aimsobjs}}

%aims! Objectives!

%%%%%%%%%%%%%%%%%%%%%%
% AIM OBJECTIVE WIBBLE
%%%%%%%%%%%%%%%%%%%%%%

This thesis has two aims. To - 

\begin{itemize}
	\item Analyse the `coordination problem' in agent-based spatial economic models
	\item Demonstrate that modelling theory impacts on how spatial economic models should be built, used and interpreted
\end{itemize}

The rationale behind each of these aims and the objectives set to achieve them, are outlined in the next two sections.

\subsection{Analyse the `coordination problem' in spatial economic models \label{sec:aim1}}

o	Objectives connected to that:
"	Present agent framework using simplest possible elements of the `standard' micro-economics framework and key agent structures; combine with spatial costs.
"	Outline and critically examine the key micro-economic and spatial economic ideas used. 

Two things:

I'm with Fowler: using ABM to develop spatial economic models is a promising way to open up new lines of research. Intuitively, it seems this would be more suited to the sort of questions we face right now. In practice, is ABM a dead end? 
Second: but (again as Fowler has found) as soon as you attempt to use ABM to model extant economic ideas, coordination problems abound. Where analytic approaches can use models that assume an equilibrium endpoint has been reached, ABM models are stuck trying to find a path to stable outcomes. 

The issues: timing structure, decision structure, how they inter-relate. That may be about it.

Solving the coordination problem for economic actors in space is difficult.
o	It presents particular challenges that non-spatial problems don't, though there may be workarounds.
o	I'm specifically trying to push that a little bit further forward. Not much, perhaps, but that's what I'm trying. 
"	I have plenty to write about as regards this: lots of modelling theory, hayek, etc. Plenty. 
o	But that does mean arguing for why this sort of modelling is about Krugman-esque "helping with thinking straight", nothing too straight.

The modelling is there as a way to examine the issues and start to get at some answers. So.

Building anything like general agent models requires understanding what happens outside of a priori equilibrium assumptions, but it doesn't require over-complicating matters (because even with the simplest assumptions, coordination problems are still legion) or declaring an economic year zero. 

The problem in a nutshell: a bunch of economic actors, with their own desires and requirements, in a decentralised system. How do they produce global economically desirable outcomes? Not necessarily optimal, just half-decent and reasonably stable?
Slightly off-topic: the more I think of it, the more it seems Fowler's approach is a) great because it discusses the issues well but b) we really need t move on from trying to compare to the core model, and begin to work out what ABM is good for. The comparison is vital - it's what I'm doing. But I'm doing it one step removed, going to some more basic economic ideas. 


%The hard part, then, is working out how to get the modelling framework to support the analysis. I've already made a start on that: partial parts of the puzzle, and noting how Fowler's tried to do the same (but coming from quite a different place.) 
%The problem: the various 60s-70s textbooks are jam-packed with a whole agent-based research agenda. I could do with picking out what I want to try.
%The odd thing: so far, there's not been any real call for ABM, though it's helped me to think through the problems. I wonder if it'll turn out to be a scaffold, a useful thinking tool?

Q: the coordination problem with agents - anything more than an intellectual exercise? Why is it important? Using ABM offers the potential of being a lot more fluid with model assumptions. Even if just at the thinking stage, this can lead to analytic ideas that would have been difficult to formulate without having gone through an ABM process. 

By concentrating on some already well-established micro-economic ideas (while tactically adding bits that ABM is good at), this allows me to keep the focus on the coordination question, and not get distracted by making more complex agents. 


To work on the problems of coordination inherent in agent-based models of spatial economics, while sticking to as many basic micro-economic assumptions as possible.
The background to that: analytic models of spatial economics, as they currently stand, are built on 'heroic assumptions' and are limited because of that. I still haven't found the best way of saying this. There are just a lot of things about spatial economics we don't know, and I think ABM could in theory be a useful tool - but there are a lot of things in the way of conclusively saying whether it is in fact useful, or just a massive distraction, and we should stick to analytics. Note brief paper of that person who's given up on ABM: one can see why you would.
"	Why "while sticking to as many basic micro-economic assumptions as possible?" Because there's a lot of good ideas, the most central of which I want to defend is the use of black boxes in utility and production.

The prime stuff: the basic micro-economics of people and firms, as basic as possible, with as simple space costs as possible, but not restricted to GE assumptions.The importance of numbers, time, price-taking, setting, speed of change. 

Hmm: perhaps the model aims are implicit in the first aim, or could be incorporated into it? Or would be an objective as part of that aim. Ah, yes. OK. 

\subsection{Demonstrate that modelling theory impacts on how spatial economic models should be built, used and interpreted \label{sec:aim2}}

o	Objectives connected to that:
"	Demonstrate that modelling theory impacts on how models should be built, used and interpreted

Develop understanding of the purpose and uses of economic models

two meanings: how much use are economic models - how effective are they? And, how should they be used?

economists have done a fair old chunk of thinking about this: more than ABM modellers, probably. My argument here is: this is isn't pointless navel gazing, it goes right to the heart of how we use models, and how we build and interpret them. This will be reflected in the way this thesis presents the building and interpretation of model `scenarios'.
%Thus, show how modelling theory / being self-aware is not pointless! 
 
%It may be that the coordination stuff is only one objective and there should be a larger aim. So what's my argument for it being the driving focus? Precisely that the model bits I have are headbutting this problem. It's the one that I don't understand, that needs fixing. 

\section{Intro wibble}


While some argue 'it is better to assume that moving goods is essentially costless' (Glaeser) because they are so low, others point out that 'our basic local needs have become dependent upon a hyper-complex, integrated, tightly-coupled global fabric of exchange' (Korowicz) 
humpty dumpty!

Glaeser and Kolhase go so far as to claim:
\begin{quote}
``These reduced costs, and the declining importance of the good-producing sector of the economy, means that in our view it is better to assume that moving goods is essentially costless than to assume that moving goods is an important component of the production process'' (\citealt{glaeser_cities_2004}, p.199).
\end{quote}

They also argue that ``there is little reason to doubt that this decline will continue'' (\citealt{glaeser_cities_2004}, p.197). This assumption of space costs' inevitably dropping is a widespread part of the literature. Lang points out that \gls{ge} models aim to:
\begin{quote}
``examine how the spatial distribution of economic activity changes as transportation costs slowly decline. The slow decline is meant to mirror the actual decline in transportation costs observed over the modern era of civilisation.'' (\citealt{lang_heterogeneous_2010} p.191)
\end{quote}

%note: orig paper for 'history' quote is krugman, venables 2005, globalisation and the inequality of nations
As Neary notes (\citealt{neary_hype_2001} p.536), this is a sweeping claim of telling the `history of the world part one', spanning ``first caravels, then steamships and railroads, then air freight...'' (\citealt{fujita_spatial_2001} p.253) This is a somewhat `Whig history' approach to thinking about space cost change, but is that a problem? One might cite recent analyses suggesting that oil prices may begin to creep back up again (see e.g. \citealt{sorrell_assessment_2009} for a thorough analysis of the issue of oil depletion and \citealt{kerschner_assessing_2009} for a rare examination of the potential economic consequences\footnote{Kerschner \emph{et al}, using an input-output model, find surprising results: agriculture is more resilient to oil cost changes than other sectors. Note, however, that no space cost or space change factors enter into the analysis.}). However, even without any specific reason to think costs may go up as well as down, thinking about space cost change `in the round' opens up interesting theoretical avenues.

\section{intro bits}

GE ``aims to explain the formation of various forms of economic agglomeration in geographical space using a general equilibrium framework'' (\citealt{fujita_frontiers_2005}, p.377). This framework is the main fracture line between GE and ABM. GE is impossible without it; ABM not only manages outside of equilibrium, many argue that, precisely because it can do this, it can more closely mirror reality. (E.g. Tesfatsion  - see chapter~\ref{chapter:ABM}.)  